\documentclass[11pt]{article}
\usepackage[a4paper,margin=27mm]{geometry}
\usepackage[T1]{fontenc}
\usepackage[utf8]{inputenc}
\usepackage{amsmath,amssymb,amsthm,booktabs,array}
\usepackage[colorlinks=true,urlcolor=blue,citecolor=blue,linkcolor=blue]{hyperref}
\setlength{\emergencystretch}{2em}
\newtheorem{theorem}{Theorem}
\newtheorem{proposition}[theorem]{Proposition}
\newtheorem{corollary}[theorem]{Corollary}
\newcommand{\A}{\mathcal A}
\newcommand{\Pp}{\mathcal P}
\title{Counterexamples to Korzekwa et al. (2019),\break
Section III.B, p.~6, Eq. (30): extending constructions of\break
Bengtsson et al. (2005) and Korzekwa et al. (2019)}
\author{Haobo Ma\\\small ChronoAI Pte Ltd, Singapore
\and Wenlin Zhang\\\small National University of Singapore\\
\small ChronoAI Pte Ltd, Singapore}
\date{6 November 2026}
\begin{document}
\maketitle
\begin{abstract}
For the product-structure conjecture of Korzekwa, Czach\'orski,
Pucha\l a and \.Zyczkowski (2019), Section III.B, Eq. (30) of
arXiv:1812.09083v3, we extend constructions of Bengtsson et al. (2005)
and Korzekwa et al. (2019, Appendix C.1) by identifying
a two-parameter family of counterexamples and giving a further counterexample
outside that family. All are common diagonals of three orthogonal pure states
and fail the proposed form under every vertex ordering. A rational example is
$(1/6,1/4,7/24,7/24)$; its total-variation distance from the union of proposed
product forms is at least $1/192$. The obstruction is an exact rank-one test
on the residual barycentric weights.
\end{abstract}

\section{The conjecture and its counterexample}
We refute the necessary product condition in Korzekwa et al. (2019),
Section III.B, p.~6, Eq. (30) of arXiv:1812.09083v3.
Let $\A_d^M$ denote the probability vectors that are common diagonals,
in a fixed orthonormal basis, of $M$ density matrices with pairwise
orthogonal supports. This is Definition 1 of \cite{KCPZ}.
For $d=4$ and $M=3$, write
\[
 f^j_k=\frac{1-\delta_{jk}}3,\qquad j,k\in\{0,1,2,3\}.
\]
The conjecture in Section III.B, Eq. (30) of \cite{KCPZ} proposes
\begin{equation}\label{eq:ansatz}
 p=abf^{\sigma(0)}+a(1-b)f^{\sigma(1)}
 +(1-a)bf^{\sigma(2)}+(1-a)(1-b)f^{\sigma(3)},
 \qquad a,b\in[0,1],
\end{equation}
as the form of points in $\A_4^3$. Here $\sigma$ is a permutation;
allowing it to depend on $p$ makes the proposed necessary condition weaker
than any fixed ordering convention.
The source's wording immediately before Eq. (30) is
``we conjecture that points belonging to $\A_4^3$ have the following
product structure''. Its Fig. 2(c) and Appendix C.3, Fig. 8 depict the
conjectured region. Our conclusion concerns this asserted necessity, not
the converse.

There is already a useful construction in the same paper:
Appendix C.1, Eqs. (C1)--(C12), proves that the segments from the centre
of the bounding tetrahedron to its edges lie in $\A_4^3$.
We show below that most points in this previously proved family already
violate Eq. (30), and give a short proof of the construction with an explicit
rational instance. The next theorem gives a different example, outside
those centre-to-edge segments. The underlying inclusion is their result.
Its radial subfamily $s=0$ already follows from
\cite[Section 4, Eqs. (35)--(40)]{BEKTZ}, released in February 2004 and
published in 2005. The two-parameter inclusion appears in the first
preprint of \cite{KCPZ} (21 December 2018), Appendix B.1,
Eqs. (B1)--(B12); the journal-associated version renumbers it as Appendix C.1.
The algorithm used for the source's numerical evidence is
described in \cite[Appendix A]{RGZ}.

\begin{theorem}\label{thm:main}
Put $r=\sqrt5$ and
\begin{equation}\label{eq:point}
 \alpha=\frac{5-r}{10},\qquad
 \beta=\frac{5-r}{20},\qquad
 \gamma=\frac{r-1}{4},\qquad
 p=(\alpha,\beta,\alpha,\gamma).
\end{equation}
Then $p\in\A_4^3$, but \eqref{eq:ansatz} fails for every permutation
$\sigma$ and every $a,b\in[0,1]$. In particular, the conjecture is
refuted (explicit counterexample), even when restricted to pure states and
to the strict interior of the bounding tetrahedron.
\end{theorem}

\begin{proof}
The inequalities $11/5<r<9/4$ follow by squaring. They imply
$0<\alpha,\beta,\gamma<1/3$, and direct summation gives
$2\alpha+\beta+\gamma=1$.
Let $z=e^{\pi i/5}$ and $t=(1+r)/2$. Consider the three column vectors
\begin{equation}\label{eq:states}
 x_0=\begin{pmatrix}\sqrt\alpha\\\sqrt\beta\\\sqrt\alpha\\\sqrt\gamma\end{pmatrix},\quad
 x_1=\begin{pmatrix}\sqrt\alpha\\\sqrt\beta z^2\\\sqrt\alpha z^4\\\sqrt\gamma z^7\end{pmatrix},\quad
 x_2=\begin{pmatrix}-\sqrt\alpha\\\sqrt\beta z^9\\\sqrt\alpha z^3\\\sqrt\gamma z^9\end{pmatrix}.
\end{equation}
All have squared coordinate moduli $p$, hence norm one. The elementary
relations $z^5=-1$, $z+\overline z=t$, and $|z|=1$ give
\begin{equation}\label{eq:powers}
 z^2=tz-1,\quad z^3=tz-t,\quad z^4=z-t,\quad
 z^7=1-tz,\quad z^9=t-z.
\end{equation}
For completeness, $2\cos(\pi/5)=t$ can also be obtained from
$z^4-z^3+z^2-z+1=0$: dividing by $z^2$ shows that
$s=z+\overline z$ satisfies $s^2-s-1=0$, and $s>0$ selects $t$.
The scalar identities needed below are
\begin{equation}\label{eq:weights}
 \alpha=2\beta,\qquad \alpha t=\beta+\gamma,\qquad
 (\gamma-\beta)t=\alpha,\qquad t^2=t+1.
\end{equation}
They follow by substituting $r^2=5$. Using \eqref{eq:powers},
\begin{align*}
 \langle x_0,x_1\rangle
 &=\alpha+\beta z^2+\alpha z^4+\gamma z^7\\
 &=(\alpha-\beta-\alpha t+\gamma)
   +(\beta t+\alpha-\gamma t)z=0,\\
 \langle x_0,x_2\rangle
 &=-\alpha+(\beta+\gamma)z^9+\alpha z^3\\
 &=\bigl[-\alpha+(\beta+\gamma-\alpha)t\bigr]
   +(\alpha t-\beta-\gamma)z=0,\\
 \langle x_1,x_2\rangle
 &=-\alpha+\beta z^7+\alpha z^9+\gamma z^2
   =-\langle x_0,x_1\rangle=0.
\end{align*}
For the second constant coefficient, \eqref{eq:weights} reduces it to
$\alpha(t^2-t-1)$. Thus the rank-one matrices
$\rho_j=x_jx_j^*$ are positive semidefinite, have trace one and common
diagonal $p$, and satisfy $\rho_i\rho_j=0$ when $i\ne j$.
This proves $p\in\A_4^3$.

To exclude product form, set $q_k=1-3p_k$. These are the unique barycentric
weights in $p=\sum_k q_k f^k$, since $\sum_kq_k=1$ and
$p_k=(1-q_k)/3$. Here
\begin{equation}\label{eq:q}
 q=(A,B,A,C),\qquad
 A=\frac{3r-5}{10},\quad B=\frac{5+3r}{20},\quad
 C=\frac{7-3r}{4}.
\end{equation}
Equation \eqref{eq:ansatz} would imply
$q_{\sigma(0)}q_{\sigma(3)}=q_{\sigma(1)}q_{\sigma(2)}$.
There are just three partitions of four coordinates into two unordered
pairs. Their product differences, in a fixed order, are
\begin{align}\label{eq:gaps}
 q_0q_1-q_2q_3&=\frac{21-9r}{10}>0,\notag\\
 q_0q_2-q_1q_3&=\frac{33-15r}{40}<0,\\
 q_0q_3-q_1q_2&=-\frac{21-9r}{10}<0.\notag
\end{align}
The signs follow from $11/5<r<9/4<7/3$. Hence no permutation can
satisfy the required equality. All $q_k$ are positive, so $p$ is strictly
inside $\Pp_4^3=\operatorname{conv}\{f^0,f^1,f^2,f^3\}$ as well.
\end{proof}

\section{What the product condition actually tests}
\begin{proposition}\label{prop:rankone}
Suppose $0\le p_k\le1/3$ and $\sum_kp_k=1$, and put $q_k=1-3p_k$.
For a fixed ordering $\sigma$, \eqref{eq:ansatz} holds if and only if
\[
 \det\begin{pmatrix}
 q_{\sigma(0)}&q_{\sigma(1)}\\
 q_{\sigma(2)}&q_{\sigma(3)}
 \end{pmatrix}=0.
\]
Consequently membership in the union of all product forms is equivalent
to at least one of the three equal-product pairings being zero.
\end{proposition}
\begin{proof}
Necessity follows from the four weights in \eqref{eq:ansatz}.
Conversely write the displayed matrix as $(u\ v;w\ h)$, with
$u+v+w+h=1$ and nonnegative entries. Set $a=u+v$, $b=u+w$.
Then $a,b\in[0,1]$ and $u-ab=uh-vw=0$. Subtracting from the row
and column sums gives $v=a(1-b)$, $w=(1-a)b$ and
$h=(1-a)(1-b)$. This includes zero rows and columns.
\end{proof}

\begin{corollary}\label{cor:distance}
Let $S$ be the union of the product forms \eqref{eq:ansatz}, over all
permutations and parameters. For the vector in \eqref{eq:point},
\[
 \inf_{p'\in S}\frac12\|p-p'\|_1\ \ge\
 \frac{15\sqrt5-33}{240}>0.
\]
\end{corollary}
\begin{proof}
The smallest absolute difference in \eqref{eq:gaps} is
$m=(15r-33)/40$: the difference between the first absolute value and
$m$ is $(117-51r)/40>0$, since $r<9/4$.
For probability vectors $q,q'$, any fixed pairing polynomial
$D(q)=q_iq_j-q_kq_l$ satisfies
$|D(q)-D(q')|\le\|q-q'\|_1$, by expanding each product difference
and bounding its coefficients by one. For $p'\in S$, one such
$D$ vanishes at $q'=1-3p'$, but has magnitude at least $m$ at $q$.
Thus $m\le\|q-q'\|_1=3\|p-p'\|_1$, as claimed.
\end{proof}

\section{Counterexamples from a previously proved family}
\begin{proposition}[Korzekwa et al. (2019), Appendix C.1]\label{prop:family}
For $0\le s\le1/2$ and $0\le t\le1/3$, put
\[
 a=\frac{1-3t}{4}+st,\qquad
 b=\frac{1-3t}{4}+(1-s)t,\qquad h=\frac{1+t}{4}.
\]
Then $p(s,t)=(a,b,h,h)\in\A_4^3$, with pure-state witnesses.
If $0\le s<1/2$ and $0<t<1/3$, it fails \eqref{eq:ansatz} for every
permutation.
\end{proposition}
\begin{proof}
The membership assertion is the inclusion proved in \cite[Appendix C.1]{KCPZ}.
Here is a short proof that also handles the endpoints.
Set $R=\sqrt{t(1+t)}/2$. The inequalities
\[
 |a-b|=t(1-2s)\le R\le a+b=\frac{1-t}{2}
\]
follow from $4t\le1+t$ and $3t\le1$. Consequently a unit complex number
$u$ can be chosen with $|a+bu|=R$: the attainable magnitudes of $a+bu$
fill $[|a-b|,a+b]$. Since $R/h\le2$, choose unit complex numbers $v,w$
with $v+w=-(a+bu)/h$. Their real inner product is
\[
 \operatorname{Re}(\overline v w)=\frac{|v+w|^2-2}{2}
 =\frac{t-1}{t+1}.
\]
The vectors with amplitude squares $(a,b,h,h)$ and phase rows
$(1,1,1,1)$, $(1,u,v,w)$, $(1,u,w,v)$ have norm one.
The first vector is orthogonal to the other two by the sum defining $v+w$;
the remaining inner product is
\[
 a+b+2h\operatorname{Re}(\overline v w)
 =\frac{1-t}{2}+\frac{t-1}{2}=0.
\]
For exclusion, write $c=(1-3t)/4$. The residual vector is
\[
 q=(c+3t(1-s),c+3ts,c,c).
\]
Its three pairing differences are
\[
 3ct+9t^2s(1-s),\qquad 3ct(1-2s),\qquad 3ct(1-2s).
\]
All are positive in the stated open parameter range. Theorem~\ref{thm:main}'s
pairing argument excludes every ordering.
\end{proof}

\begin{corollary}\label{cor:rational}
The rational diagonal $p_*=(1/6,1/4,7/24,7/24)$ belongs to $\A_4^3$
but has no product representation under any ordering. Its distance from
the union $S$ of all forms \eqref{eq:ansatz} satisfies
\[
 \inf_{p'\in S}\tfrac12\|p_*-p'\|_1\ge\frac1{192}.
\]
\end{corollary}
\begin{proof}
Take $s=1/4,t=1/6$ in Proposition~\ref{prop:family}.
For an explicit witness, set
\[
 \omega=\frac{-1+i\sqrt3}{2},\quad
 v=-\frac{(1+3i\sqrt3)(1+i\sqrt6)}{14},\quad
 w=-\frac{(1+3i\sqrt3)(1-i\sqrt6)}{14}.
\]
Then $|\omega|=|v|=|w|=1$, $v+w=-(1+3i\sqrt3)/7$ and
$\operatorname{Re}(\overline v w)=-5/7$.
Use the three phase rows in the preceding proof with
$(a,b,h)=(1/6,1/4,7/24)$ and $u=\omega$.
The residual vector is $(1/2,1/4,1/8,1/8)$, with pairing differences
$7/64,1/32,1/32$. The Lipschitz argument of
Corollary~\ref{cor:distance}, now with $m=1/32$, gives the bound.
\end{proof}

The example in Theorem~\ref{thm:main} is not a permutation of a member
of the centre-to-edge family: in that family the repeated coordinates $h,h$
are maximal, whereas its repeated coordinates $\alpha,\alpha$ are strictly
between $\beta$ and $\gamma$.

For $s=0$ the three residual pairing differences all equal
$3ct>0$. Thus the radial subfamily already supplies counterexamples
using the construction of Bengtsson et al. (2004/2005).
The product-pairing equations also describe boundaries of a different
set, the semigroup-compatible Pauli channels in
\cite[Section IV, Eq. (23)]{PRZ}; that use does not characterize $\A_4^3$.

\begin{corollary}\label{cor:unitary}
For either explicit diagonal above, $T(p)=(p\ p\ p\ 1-3p)$ is
unistochastic. The largest number of perfectly distinguishable states with
that diagonal is exactly three. Every stochastic matrix having such a
column admits at least three perfectly distinguishable channels without
entangled inputs. For $T(p)$ that number is four.
\end{corollary}
\begin{proof}
Complete the three orthonormal vectors to a unitary $U$. Row normalization
gives $|U_{k3}|^2=1-3p_k$, as in \cite[Lemma 5]{KCPZ}.
Four orthogonal-support density matrices in dimension four would all be
rank one and sum to the identity, forcing $4p_k=1$ for all $k$, which neither
diagonal satisfies.
For a stochastic matrix $T$ with the selected column $l$, choose unit vectors
$y_j^{(n)}$ with amplitude squares $T_{\cdot j}$, using the three orthogonal
witnesses when $j=l$. The Kraus operators
$K_j^{(n)}=|y_j^{(n)}\rangle\langle j|$ define trace-preserving channels
with classical action $T$ and orthogonal outputs on $|l\rangle$.
This is \cite[Proposition 19]{KCPZ} specialized to these diagonals.
Finally, for the unistochastic $T(p)=|U|^2$, the four unitaries $D_jU$, where
$D_j=\operatorname{diag}(1,i^j,i^{2j},i^{3j})$, have the same action and
send $U^*(1,1,1,1)^\top/2$ to four orthogonal Fourier vectors.
No unassisted protocol has more than four orthogonal outputs in dimension
four. This is the construction of \cite[Proposition 15]{KCPZ}.
\end{proof}

\begin{corollary}[Korzekwa et al. (2019), Appendix C.2]\label{cor:face}
No point in the relative interior of a two-dimensional face of $\Pp_4^3$
belongs to $\A_4^3$.
\end{corollary}
\begin{proof}
After permutation such a point has $p_0=1/3$ and
$0<p_j<1/3$ for $j=1,2,3$. Lemma 6 of \cite{KCPZ} reduces the question
to pure states. Make their first coordinates $1/\sqrt3$ and their first
column real by phase choices. The truncated three columns have Gram matrix
with diagonal $2/3$ and off-diagonal $-1/3$, so their sum is zero.
In each remaining row the three entries therefore have equal nonzero modulus
and sum zero; its phases are $(1,\omega,\omega^2)$ or
$(1,\omega^2,\omega)$ with $\omega=e^{2\pi i/3}$.
The imaginary part of the first two columns' inner product requires the
three positive weights $p_1,p_2,p_3$ to split into two sets of equal weight
$1/3$. One of these sets has at most one element, contradicting $p_j<1/3$.
This is a short proof of the source's face exclusion, not a new exclusion.
\end{proof}

\section{Pure witnesses and the excluded radial segment}
\begin{proposition}\label{prop:pure}
Every $p\in\A_4^3$ has three orthogonal pure-state witnesses.
Consequently $p\in\A_4^3$ if and only if
$B(p)=(p\ p\ p\ \mathbf1-3p)$ is unistochastic.
\end{proposition}
\begin{proof}
The ranks of three nonzero density matrices with orthogonal supports
sum to at most four. Unless all are pure, their ranks are $(2,1,1)$.
Let $x,y$ be unit vectors for the two pure states, so
$|x_k|^2=|y_k|^2=p_k$. A diagonal unitary can make $y=\overline x$:
at each nonzero coordinate choose its phase $d_k$ with
$d_k^2=\overline{x_k}/y_k$, and choose it arbitrarily when $p_k=0$.
This transformation preserves all diagonals and orthogonality.
The two-dimensional complement $S$ of $\operatorname{span}\{x,\overline x\}$
is invariant under conjugation. If $\rho$ is the remaining state, then
$R=(\rho+\overline\rho)/2$ is real, positive semidefinite, trace one,
supported on $S$ and has diagonal $p$.
Real spectral decomposition on $S$ gives
$R=\lambda aa^{\mathsf T}+(1-\lambda)bb^{\mathsf T}$,
where $a,b\in S$ are real orthonormal vectors and $0\le\lambda\le1$.
Then $z=\sqrt\lambda\,a+i\sqrt{1-\lambda}\,b$ is a unit vector in $S$
and $|z_k|^2=R_{kk}=p_k$. Replacing $\rho$ by $zz^*$ gives three pure
witnesses. Completing their columns to a unitary gives $B(p)$;
conversely its first three unitary columns give the required witnesses.
\end{proof}

\begin{corollary}\label{cor:radial}
For $-1/9\le t<0$, the probability vector
$p_t=(1-3t,1+t,1+t,1+t)/4$ does not belong to $\A_4^3$.
The endpoint $p_0=(1,1,1,1)/4$ does belong to $\A_4^3$.
\end{corollary}
\begin{proof}
For a pure triple at $p_t$, dephasing one column and the first row
reduces the remaining orthogonality conditions to the phase system
in \cite[Section 4, Eqs. (35)--(38)]{BEKTZ}, with
$L=(1-3t)/(1+t)>1$. The lemma in its Appendix C proves that this
system has no solution for $L>1$. Proposition~\ref{prop:pure} extends
the exclusion from pure witnesses to arbitrary density matrices.
At $t=0$, three columns of the Fourier matrix give pure witnesses.
Thus the exclusion interval printed after Eq. (C37) in \cite{KCPZ}
must omit its central endpoint.
\end{proof}

\section{Coverage and further questions}
The counterexample refutes necessity of Eq. (30), including the full-region
interpretation of its plots. It does not refute the independently proved
inclusions or exclusions in Appendix C.1--2 of \cite{KCPZ}.
The relation to unistochastic matrices places the question in the
dimension-four geometry studied by Bengtsson, Ericsson, Ku\'s, Tadej and
\.Zyczkowski \cite{BEKTZ}; our obstruction is the explicit failure of the
residual-weight rank-one condition, rather than a failure of unistochasticity.
The following table separates the asserted conjecture, previously proved
results and statements not posed by the source. No replacement description
of the full region is asserted.
\begin{center}
\small
\begin{tabular}{@{}p{.48\linewidth}p{.46\linewidth}@{}}
\toprule
Source item in \cite{KCPZ} & Coverage\\
\midrule
Section III.B, Eq. (30); Fig. 2(c); Appendix C.3, Fig. 8
 & Refuted (explicit counterexamples), Theorem~\ref{thm:main} and
 Proposition~\ref{prop:family}.\\
Necessity under point-dependent ordering; pure/interior restrictions
 & Not asserted separately in the source; fails for both explicit diagonals.\\
Proposition 4: $\A_4^3\subseteq\Pp_4^3$
 & Preserved; all four coordinates satisfy $p_k<1/3$.\\
Appendix C.1, Eqs. (C1)--(C12): centre-to-edge inclusion
 & Reproved; prior proof Korzekwa et al. (2019).
 Its generic points also fail Eq. (30).\\
Appendix C.2: proper-face exclusion; radial exclusion
 & Face: reproved; prior proof Korzekwa et al. (2019).
 Radial: prior pure-state proof Bengtsson et al. (2004/2005) cited;
 mixed states follow from Proposition~\ref{prop:pure}.\\
Appendix C.2, interval following Eqs. (C36)--(C37), inclusion of $t=0$
 & Refuted-and-corrected at this endpoint: $\eta$ is included;
 the exclusion uses $t<0$ (Corollary~\ref{cor:radial}).\\
Lemma 5; Propositions 15 and 19
 & Prior proofs cited; explicit instances in Corollary~\ref{cor:unitary}.\\
Proposition 12: coarse-graining implication
 & Prior proof Korzekwa et al. (2019); unchanged, not reproved.\\
Proposition 8; Proposition 13 and Remark 14
 & Prior results cited; no new global bound or centre-neighbourhood result.\\
Conjecture 9, Eq. (31); Remark 10: star-shapedness
 & Not covered by this construction.\\
Sufficiency of Eq. (30); a full description of $\A_4^3$
 & Not established here.\\
Rank-one equivalence; distance estimates; exact state capacity;
 pure witnesses throughout $\A_4^3$
 & Not posed; proved here.\\
Other Section IV results; Section V questions on noise and resources
 & Not covered beyond Corollary~\ref{cor:unitary}.\\
\bottomrule
\end{tabular}
\end{center}
For the face argument after Eq. (C34) in \cite{KCPZ}, the normalized
positive weights satisfy $\sum_j a_j=1$ and $\sum_j a_j\epsilon_j=0$
with $\epsilon_j\in\{-1,1\}$. The two sign classes each have weight
$1/2$, and one class is a singleton. Thus one $a_j=1/2$; the other two
need only sum to $1/2$. Their equality is unnecessary for the exclusion.

Proposition~\ref{prop:pure} connects Conjecture 9 in dimension four to
the special repeated-column matrices $B(p)$. In particular,
$B(\lambda p+(1-\lambda)\eta)=\lambda B(p)+(1-\lambda)W_4$.
The star-shapedness theorems of \cite[Theorem 14, Corollary 16,
Theorem 17]{RMKPZ} concern bracelet and factorisable matrices,
dimension-three unistochastic matrices, and dimension-four circulant
unistochastic matrices, respectively. None establishes star-shapedness
of this repeated-column subset. The distinction between the full
unistochastic set and its circulant subset is also made in
\cite[Section 4.3]{Salazar}. Conjecture 9 remains outside the conclusions
of this note.

\paragraph{Mathematical verification.}
The declaration \texttt{result} in the pinned
\href{https://github.com/the-omega-institute/trureturing/blob/73ec20a01fc35e0f53bc258c605821728f3c33b6/D5/S3/Quantum/Measurements/ThreeDistinguishabilityProductRefutation.lean}{Lean source}
states $\neg\exists\sigma\,\forall p\in\texttt{InRegion},\,
\texttt{ProductForm}(\sigma,p)$. Its definition of \texttt{InRegion}
uses three positive-semidefinite trace-one matrices with pairwise zero products
and the specified real diagonals. Its proof constructs the diagonal
\eqref{eq:point}; the private lemma \texttt{point\_not\_product} excludes
every permutation for that same point. Thus the literal public declaration
has a fixed outer ordering, while the witness-specific proof gives the stronger
pointwise obstruction. The rank-one criterion, distance estimates, family
and rational example, pure-state reduction, radial exclusion,
exact state capacity and channel instances have written
proofs here; they are not additional statements of \texttt{result}.
The accompanying \texttt{verify.py} checks the weighted Gram matrix,
the three pairing differences and all 24 permutations for each explicit
diagonal in exact number-field arithmetic. It also checks the scalar identities,
strict sign bounds and the family pairing polynomials.
The geometric statements follow from the written arguments,
not from numerical sampling.

\begin{thebibliography}{6}
\bibitem{KCPZ}
K. Korzekwa, S. Czach\'orski, Z. Pucha\l a and K. \.Zyczkowski,
\emph{Distinguishing classically indistinguishable states and channels},
J. Phys. A: Math. Theor. \textbf{52} (2019), 475303.
\href{https://doi.org/10.1088/1751-8121/ab30f7}{doi:10.1088/1751-8121/ab30f7}.
\href{https://arxiv.org/abs/1812.09083v3}{arXiv:1812.09083v3},
Section III.B, Eq. (30); Appendix C.3.
First preprint: \href{https://arxiv.org/abs/1812.09083v1}{arXiv:1812.09083v1}
(21 December 2018), Appendix B.1.
\bibitem{BEKTZ}
I. Bengtsson, \AA. Ericsson, M. Ku\'s, W. Tadej and K. \.Zyczkowski,
\emph{Birkhoff's polytope and unistochastic matrices, $N=3$ and $N=4$},
Commun. Math. Phys. \textbf{259} (2005), 307--324.
\href{https://doi.org/10.1007/s00220-005-1392-8}{doi:10.1007/s00220-005-1392-8}.
\href{https://arxiv.org/abs/math/0402325v2}{arXiv:math/0402325v2},
Section 1, Eqs. (1)--(2), and Section 4.
Preprint released 19 February 2004; version 2 dated 24 February 2004.
\bibitem{RGZ}
G. Rajchel, A. G\k{a}siorowski and K. \.Zyczkowski,
\emph{Robust Hadamard matrices, unistochastic rays in Birkhoff polytope
and equi-entangled bases in composite spaces},
Math. Comput. Sci. \textbf{12} (2018), 473--490.
\href{https://doi.org/10.1007/s11786-018-0384-y}{doi:10.1007/s11786-018-0384-y}.
\href{https://arxiv.org/abs/1804.10715}{arXiv:1804.10715}, Appendix A.
\bibitem{PRZ}
Z. Pucha\l a, \L. Rudnicki and K. \.Zyczkowski,
\emph{Pauli semigroups and unistochastic quantum channels},
Phys. Lett. A \textbf{383} (2019), 2376--2381.
\href{https://doi.org/10.1016/j.physleta.2019.04.057}{doi:10.1016/j.physleta.2019.04.057}.
\href{https://arxiv.org/abs/1903.12448v1}{arXiv:1903.12448v1}, Section IV, Eq. (23).
\bibitem{RMKPZ}
G. Rajchel-Mieldzio\'c, K. Korzekwa, Z. Pucha\l a and K. \.Zyczkowski,
\emph{Algebraic and geometric structures inside the Birkhoff polytope},
J. Math. Phys. \textbf{63} (2022), 012202.
\href{https://doi.org/10.1063/5.0046581}{doi:10.1063/5.0046581}.
\href{https://arxiv.org/abs/2101.11288v1}{arXiv:2101.11288v1},
Theorem 14, Corollary 16 and Theorem 17.
\bibitem{Salazar}
R. Salazar, J. Czartowski, R. Ravell Rodr\'iguez, G. Rajchel-Mieldzio\'c,
P. Horodecki and K. \.Zyczkowski,
\emph{Quantum resource theories beyond convexity},
Quantum \textbf{10} (2026), 2104.
\href{https://doi.org/10.22331/q-2026-05-13-2104}{doi:10.22331/q-2026-05-13-2104}.
\href{https://arxiv.org/abs/2405.05785v6}{arXiv:2405.05785v6}, Section 4.3.
\end{thebibliography}
\paragraph{License.} CC BY 4.0 (Creative Commons Attribution 4.0 International).
\end{document}
